\subsection{Literaturrecherche}
Welche aktuellen Ansätze gibt es 
\begin{itemize}
    \item zur Generierung von synthetischen klinischen Daten (v.a. große Sprachmodelle und Agentensysteme, sowohl klassisch als auch LLM-basiert)?
    \item für Schnittstellen zum Export von medizinisch relevanten Daten und zur Integration von Sensordaten in Krankenhausinformationssystem? 
\end{itemize}

\subsection{Definition der Qualitätskriterien für Daten in dem Heidelberger Lehre-KIS}
Welche Anforderungen müssen die Sensordaten und synthetisch generierte Daten im Heidelberger Lehre-KIS erfüllen, damit ein didaktisches Konzept für die Übungsaufgaben entwickelt werden kann?

\subsection{Prototypische Implementierung und Erzeugung von Beispiel-Datensätzen}
Es sollen synthetische Daten über Sprachmodelle und nicht-Sprachmodell-basierte Agenten erzeugt werden.

Es sollen konkrete Sensordaten aus medizintechnischen Geräten oder Wearables über Schnittstellen in das Lehre-KIS eingelesen werden.

\subsection{Bewertung der erprobten Ansätze}
Verwendung der definierten Qualitätskriterien um die implementierten Ansätze zu bewerten und das am besten geeignete Verfahren für das Heidelberger Lehre-KIS vorschlagen. Zusätzlich sollen auch Aspekte wie Energieverbrauch und Komplexität der Anwendung gegeneinander abgewogen werden.